\DocumentMetadata{
  pdfversion=2.0,
  pdfstandard=ua-2,
  lang=en-US,
  tagging=on,
  tagging-setup={math/setup=mathml-SE,tabsorder=structure}
}
\documentclass[11pt,letterpaper]{article}
\usepackage[margin=0.68in,includefoot]{geometry}
\usepackage{newcomputermodern}
\usepackage{amsmath,amscd,mathtools}
\usepackage{tikz,adjustbox}
\providecommand{\square}{\mdlgwhtsquare}
\usepackage[hidelinks]{hyperref}
\hypersetup{
  pdftitle={Logic PhD Exam, May 1993},
  pdfauthor={University of Florida Department of Mathematics},
  pdfsubject={Graduate mathematics examination},
  pdfdisplaydoctitle=true,
  bookmarksopen=true
}
\setcounter{secnumdepth}{0}
\setlength{\parindent}{0pt}
\setlength{\parskip}{0.28em}
\setlength{\emergencystretch}{3em}
\setlength{\leftmargini}{2.15em}
\setlength{\leftmarginii}{2.65em}
\setlength{\leftmarginiii}{3em}
\renewcommand{\labelenumi}{\textbf{\theenumi.}}
\pagestyle{empty}
\raggedbottom
\allowdisplaybreaks
\newenvironment{examproblems}[1]{%
  \begin{enumerate}%
  \setcounter{enumi}{#1}%
  \setlength{\topsep}{0.35em}%
  \setlength{\partopsep}{0pt}%
  \setlength{\itemsep}{0.48em}%
  \setlength{\parsep}{0.18em}%
}{\end{enumerate}}
\newenvironment{examsubparts}{%
  \begin{enumerate}%
  \setlength{\topsep}{0.18em}%
  \setlength{\partopsep}{0pt}%
  \setlength{\itemsep}{0.16em}%
  \setlength{\parsep}{0.1em}%
}{\end{enumerate}}
\newenvironment{examinstructions}{%
  \par\begingroup\itshape
}{%
  \par\endgroup\vspace{0.18em}
}
\makeatletter
\renewcommand\section{\@startsection{section}{1}{0pt}%
  {-1.25ex plus -.25ex minus -.1ex}%
  {0.55ex plus .1ex}%
  {\normalfont\large\bfseries}}
\renewcommand{\@maketitle}{%
  \newpage\null\vskip -1em%
  {\footnotesize\bfseries
    \href{https://www.ufl.edu/}{University of Florida}, \href{https://math.ufl.edu/}{Department of Mathematics}\par}%
  \vskip -0.5em%
  \tagstructbegin{tag=Title}%
    {\LARGE\bfseries\href{https://gma.math.ufl.edu/exams/logic-phd/}{Logic PhD Exam}, May 1993\par}%
  \tagstructend%
  \vskip 0.25em%
  \tagmcbegin{artifact}%
    \hrule height 1pt%
  \tagmcend%
  \vskip 1em%
}
\makeatother
\title{Logic PhD Exam, May 1993}
\author{}
\date{}
\begin{document}
\maketitle
\thispagestyle{empty}
\begin{examinstructions}
This examination is divided into 4 segments, with 3 questions in each segment. Answer any 2 questions from each segment.
\end{examinstructions}
\section{1. General Logic}
\begin{examproblems}{0}
\item
State Gödel’s incompleteness theorem. Explain the importance of Gödel numbering.
\item
State and prove the compactness theorem for Predicate logic. You may assume the completeness theorem.
\item
Prove that if~\(T\) is a countable theory with arbitrarily large finite models then~\(T\) has an infinite model.
\end{examproblems}
\section{2. Model Theory}
\begin{examproblems}{0}
\item
State and sketch the proof of the Łoś theorem on ultraproducts.
\item
State and sketch a proof of the elementary substructure version of the downward Löwenheim–Skolem Theorem.
\item
Prove that the theory of dense linear orders without endpoints is complete.
\end{examproblems}
\section{3. Set Theory}
\begin{examproblems}{0}
\item
Prove that the well ordering principle implies the axiom of choice.
\item
Prove, using the axioms of ZFC, that if~\(X\) is any set then the transitive closure of~\(X\) is also a set.
\item
Show that a countable union of countable sets is countable, pointing out any use of the axiom of choice.
\end{examproblems}
\section{4. Recursion Theory}
\begin{examproblems}{0}
\item
Show that a function \(\phi:\omega\longrightarrow\omega\) is (partial) recursive iff its graph is recursively enumerable, and that if~\(\phi\) is total and recursive then its graph is recursive.
\item
State the halting problem, and show that it is undecidable. You may use the fact that there is a recursive enumeration of the (partial) recursive functions.
\item
Show that if \(A\ne\varnothing\) is the domain of a partial recursive function~\(f\) then there is a total recursive function~\(g\) such that~\(A\) is the range of~\(g\).
\end{examproblems}
\end{document}
